% --- Archon physics formalization source begin ---
% archon:physics
% archon:covers PhyXMiniProblems/problem_phyx_mini_0196.lean
% archon:source-report reports/phyx_mini/problem_phyx_mini_0196.source.json

\chapter{Physics problem phyx\_mini\_0196}
\label{ch:PhyXMiniProblems_problem_phyx_mini_0196}

\paragraph{Problem source.}
\#\# Physical scenario

The siren is moving away from the listener with a speed of 45 \textbackslash{}, \textbackslash{}mathrm\{m/s\} relative to the air, and the listener is moving toward the siren with a speed of 15 \textbackslash{}, \textbackslash{}mathrm\{m/s\} relative to the air.

\#\# Question

What frequency does the listener hear?

\#\# Answer choices

- A: A: 267Hz

- B: B: 277Hz

- C: C: 274Hz

- D: D: 268Hz

\#\# Figure caption (auxiliary; use the image as primary evidence)

The image depicts two cars on a straight road, illustrating a scenario often used in physics to explain the Doppler effect.

- On the left, there is a car labeled "Listener" (L). It is moving to the right with a velocity of \textbackslash{}( v\_L = 15 \textbackslash{}, \textbackslash{}text\{m/s\} \textbackslash{}). The frequency perceived by the listener is labeled as \textbackslash{}( f\_L = ? \textbackslash{}).

- To the right, there is a "Police car" (S) also moving rightward with a velocity of \textbackslash{}( v\_S = 45 \textbackslash{}, \textbackslash{}text\{m/s\} \textbackslash{}). 

- There is an arrow labeled "+" indicating the direction from the listener (L) to the source (S).

The cars are shown with motion lines to emphasize their movement, and there is an explicit horizontal line that the vehicles sit on, suggesting the road.

\#\# Dataset metadata

Wave Properties; Physical Model Grounding Reasoning, Multi-Formula Reasoning

\paragraph{Recorded answer/context.}
B: B: 277Hz

\paragraph{Figure/image path.}
/root/proposal\_for\_physic/science-mango/phyx\_mini\_run/phyx\_data/test\_image/196.png

\paragraph{Formalization target.}
create a compiling Lean file with sorry bodies at `PhyXMiniProblems/problem\_phyx\_mini\_0196.lean`. The Lean declarations must preserve the physical quantities, dimensions or dimensional roles, figure labels, governing-law hypotheses, and final relation expressed by this problem.
Use Mathlib/Physlib names found through LeanExplore where available. If a physics API is missing, introduce faithful local abstractions rather than scalar placeholder aliases.

\begin{theorem}[Physics formalization target]
\label{thm:physics:phyx_mini_0196:target}
\lean{PhyXMiniProblems.ProblemPhyXMini0196.heardFrequency_matches_recordedChoiceB}
\uses{def:physics:phyx-mini-0196:phyxminiproblems-problemphyxmini0196-sirenlistenerdopplersetup, def:physics:phyx-mini-0196:phyxminiproblems-problemphyxmini0196-matchesprimaryfigure, def:physics:phyx-mini-0196:phyxminiproblems-problemphyxmini0196-usespriorsirencalibration, def:physics:phyx-mini-0196:phyxminiproblems-problemphyxmini0196-hassubsonicphysicalparameters, def:physics:phyx-mini-0196:phyxminiproblems-problemphyxmini0196-satisfiesleftwardacousticdopplerlaw, def:physics:phyx-mini-0196:phyxminiproblems-problemphyxmini0196-recordeddatasetanswer, def:physics:phyx-mini-0196:phyxminiproblems-problemphyxmini0196-matchesanswerchoice}
The assigned autoformalize agent should translate this physics problem into Lean declarations in the covered file, with theorem and lemma proof bodies written as `by sorry`.
\end{theorem}
\begin{proof}
This is an autoformalization task, not a proof task. Produce statements that can later be proved by the physics prover without weakening the physical model.
\end{proof}

% --- Archon declaration topology begin ---
\section{Declaration topology}
\begin{definition}[Acoustic Frequency]
  \label{def:physics:phyx-mini-0196:phyxminiproblems-problemphyxmini0196-acousticfrequency}
  \lean{PhyXMiniProblems.ProblemPhyXMini0196.AcousticFrequency}
  A physical acoustic frequency, carrying inverse-time dimension
\end{definition}
\begin{proof}
  This is the declared model definition of Acoustic Frequency.
\end{proof}
\begin{definition}[Axial Position]
  \label{def:physics:phyx-mini-0196:phyxminiproblems-problemphyxmini0196-axialposition}
  \lean{PhyXMiniProblems.ProblemPhyXMini0196.AxialPosition}
  A signed physical coordinate along the horizontal road axis
\end{definition}
\begin{proof}
  This is the declared model definition of Axial Position.
\end{proof}
\begin{definition}[Axial Velocity]
  \label{def:physics:phyx-mini-0196:phyxminiproblems-problemphyxmini0196-axialvelocity}
  \lean{PhyXMiniProblems.ProblemPhyXMini0196.AxialVelocity}
  A signed physical velocity along the road axis. A signed quantity is used instead of Physlib's unsigned DimSpeed so that the figure's positive direction and the participants' directions of motion remain explicit
\end{definition}
\begin{proof}
  This is the declared model definition of Axial Velocity.
\end{proof}
\begin{definition}[frequency Readout]
  \label{def:physics:phyx-mini-0196:phyxminiproblems-problemphyxmini0196-frequencyreadout}
  \lean{PhyXMiniProblems.ProblemPhyXMini0196.frequencyReadout}
  \uses{def:physics:phyx-mini-0196:phyxminiproblems-problemphyxmini0196-acousticfrequency}
  Read a physical acoustic frequency in inverse units of a chosen time unit
\end{definition}
\begin{proof}
  This is the declared model definition of frequency Readout.
\end{proof}
\begin{definition}[frequency In Hertz]
  \label{def:physics:phyx-mini-0196:phyxminiproblems-problemphyxmini0196-frequencyinhertz}
  \lean{PhyXMiniProblems.ProblemPhyXMini0196.frequencyInHertz}
  \uses{def:physics:phyx-mini-0196:phyxminiproblems-problemphyxmini0196-acousticfrequency, def:physics:phyx-mini-0196:phyxminiproblems-problemphyxmini0196-frequencyreadout}
  Hertz readout of an acoustic frequency
\end{definition}
\begin{proof}
  This is the declared model definition of frequency In Hertz.
\end{proof}
\begin{definition}[position Readout]
  \label{def:physics:phyx-mini-0196:phyxminiproblems-problemphyxmini0196-positionreadout}
  \lean{PhyXMiniProblems.ProblemPhyXMini0196.positionReadout}
  \uses{def:physics:phyx-mini-0196:phyxminiproblems-problemphyxmini0196-axialposition}
  Read a signed physical position in a chosen length unit
\end{definition}
\begin{proof}
  This is the declared model definition of position Readout.
\end{proof}
\begin{definition}[position In Meters]
  \label{def:physics:phyx-mini-0196:phyxminiproblems-problemphyxmini0196-positioninmeters}
  \lean{PhyXMiniProblems.ProblemPhyXMini0196.positionInMeters}
  \uses{def:physics:phyx-mini-0196:phyxminiproblems-problemphyxmini0196-axialposition, def:physics:phyx-mini-0196:phyxminiproblems-problemphyxmini0196-positionreadout}
  Metre readout of a signed axial position
\end{definition}
\begin{proof}
  This is the declared model definition of position In Meters.
\end{proof}
\begin{definition}[axial Velocity Readout]
  \label{def:physics:phyx-mini-0196:phyxminiproblems-problemphyxmini0196-axialvelocityreadout}
  \lean{PhyXMiniProblems.ProblemPhyXMini0196.axialVelocityReadout}
  \uses{def:physics:phyx-mini-0196:phyxminiproblems-problemphyxmini0196-axialvelocity}
  Read a signed axial velocity in selected length and time units
\end{definition}
\begin{proof}
  This is the declared model definition of axial Velocity Readout.
\end{proof}
\begin{definition}[axial Velocity In Meters Per Second]
  \label{def:physics:phyx-mini-0196:phyxminiproblems-problemphyxmini0196-axialvelocityinmeterspersecond}
  \lean{PhyXMiniProblems.ProblemPhyXMini0196.axialVelocityInMetersPerSecond}
  \uses{def:physics:phyx-mini-0196:phyxminiproblems-problemphyxmini0196-axialvelocity, def:physics:phyx-mini-0196:phyxminiproblems-problemphyxmini0196-axialvelocityreadout}
  Metres-per-second readout of a signed air-relative axial velocity
\end{definition}
\begin{proof}
  This is the declared model definition of axial Velocity In Meters Per Second.
\end{proof}
\begin{definition}[speed In Meters Per Second]
  \label{def:physics:phyx-mini-0196:phyxminiproblems-problemphyxmini0196-speedinmeterspersecond}
  \lean{PhyXMiniProblems.ProblemPhyXMini0196.speedInMetersPerSecond}
  Read Physlib's unsigned physical speed in metres per second
\end{definition}
\begin{proof}
  This is the declared model definition of speed In Meters Per Second.
\end{proof}
\begin{definition}[Participant]
  \label{def:physics:phyx-mini-0196:phyxminiproblems-problemphyxmini0196-participant}
  \lean{PhyXMiniProblems.ProblemPhyXMini0196.Participant}
  The two labeled participants in the primary figure
\end{definition}
\begin{proof}
  This is the declared model definition of Participant.
\end{proof}
\begin{definition}[Road Direction]
  \label{def:physics:phyx-mini-0196:phyxminiproblems-problemphyxmini0196-roaddirection}
  \lean{PhyXMiniProblems.ProblemPhyXMini0196.RoadDirection}
  The two orientations along the depicted one-dimensional road
\end{definition}
\begin{proof}
  This is the declared model definition of Road Direction.
\end{proof}
\begin{definition}[Waveform Kind]
  \label{def:physics:phyx-mini-0196:phyxminiproblems-problemphyxmini0196-waveformkind}
  \lean{PhyXMiniProblems.ProblemPhyXMini0196.WaveformKind}
  The waveform classification inherited from the siren exercise context
\end{definition}
\begin{proof}
  This is the declared model definition of Waveform Kind.
\end{proof}
\begin{definition}[Siren Listener Doppler Setup]
  \label{def:physics:phyx-mini-0196:phyxminiproblems-problemphyxmini0196-sirenlistenerdopplersetup}
  \lean{PhyXMiniProblems.ProblemPhyXMini0196.SirenListenerDopplerSetup}
  \uses{def:physics:phyx-mini-0196:phyxminiproblems-problemphyxmini0196-acousticfrequency, def:physics:phyx-mini-0196:phyxminiproblems-problemphyxmini0196-axialposition, def:physics:phyx-mini-0196:phyxminiproblems-problemphyxmini0196-axialvelocity, def:physics:phyx-mini-0196:phyxminiproblems-problemphyxmini0196-participant, def:physics:phyx-mini-0196:phyxminiproblems-problemphyxmini0196-roaddirection, def:physics:phyx-mini-0196:phyxminiproblems-problemphyxmini0196-waveformkind}
  The physical apparatus and observables needed for the one-dimensional acoustic Doppler model. heardFrequency is an unknown physical observable; the structure does not assign it a numerical value or answer choice
\end{definition}
\begin{proof}
  This is the declared model definition of Siren Listener Doppler Setup.
\end{proof}
\begin{definition}[Matches Primary Figure]
  \label{def:physics:phyx-mini-0196:phyxminiproblems-problemphyxmini0196-matchesprimaryfigure}
  \lean{PhyXMiniProblems.ProblemPhyXMini0196.MatchesPrimaryFigure}
  \uses{def:physics:phyx-mini-0196:phyxminiproblems-problemphyxmini0196-positioninmeters, def:physics:phyx-mini-0196:phyxminiproblems-problemphyxmini0196-axialvelocityinmeterspersecond, def:physics:phyx-mini-0196:phyxminiproblems-problemphyxmini0196-sirenlistenerdopplersetup}
  Primary-image geometry and readouts. The listener L lies to the left of source S; the arrow from L to S is positive; the sound reaching L travels in the reverse direction; and the air-relative velocities are v\_L = 15 m/s and v\_S = 45 m/s. No readout of the requested heard frequency occurs here
\end{definition}
\begin{proof}
  This is the declared model definition of Matches Primary Figure.
\end{proof}
\begin{definition}[Uses Prior Siren Calibration]
  \label{def:physics:phyx-mini-0196:phyxminiproblems-problemphyxmini0196-usespriorsirencalibration}
  \lean{PhyXMiniProblems.ProblemPhyXMini0196.UsesPriorSirenCalibration}
  \uses{def:physics:phyx-mini-0196:phyxminiproblems-problemphyxmini0196-frequencyinhertz, def:physics:phyx-mini-0196:phyxminiproblems-problemphyxmini0196-speedinmeterspersecond, def:physics:phyx-mini-0196:phyxminiproblems-problemphyxmini0196-sirenlistenerdopplersetup}
  The shared calibration from source panel 193 of the same police-siren sequence: the siren emits a 300 Hz sinusoid and sound travels through the air at 340 m/s. These data are separated from the current panel's figure readouts because they are not reprinted in the assigned image. This premise contains no heard-frequency or answer-choice conclusion
\end{definition}
\begin{proof}
  This is the declared model definition of Uses Prior Siren Calibration.
\end{proof}
\begin{definition}[Has Subsonic Physical Parameters]
  \label{def:physics:phyx-mini-0196:phyxminiproblems-problemphyxmini0196-hassubsonicphysicalparameters}
  \lean{PhyXMiniProblems.ProblemPhyXMini0196.HasSubsonicPhysicalParameters}
  \uses{def:physics:phyx-mini-0196:phyxminiproblems-problemphyxmini0196-frequencyinhertz, def:physics:phyx-mini-0196:phyxminiproblems-problemphyxmini0196-axialvelocityinmeterspersecond, def:physics:phyx-mini-0196:phyxminiproblems-problemphyxmini0196-speedinmeterspersecond, def:physics:phyx-mini-0196:phyxminiproblems-problemphyxmini0196-sirenlistenerdopplersetup}
  Positivity and subsonic hypotheses selecting the ordinary acoustic branch and ensuring the Doppler denominator is nonzero
\end{definition}
\begin{proof}
  This is the declared model definition of Has Subsonic Physical Parameters.
\end{proof}
\begin{definition}[Satisfies Leftward Acoustic Doppler Law]
  \label{def:physics:phyx-mini-0196:phyxminiproblems-problemphyxmini0196-satisfiesleftwardacousticdopplerlaw}
  \lean{PhyXMiniProblems.ProblemPhyXMini0196.SatisfiesLeftwardAcousticDopplerLaw}
  \uses{def:physics:phyx-mini-0196:phyxminiproblems-problemphyxmini0196-frequencyinhertz, def:physics:phyx-mini-0196:phyxminiproblems-problemphyxmini0196-axialvelocityinmeterspersecond, def:physics:phyx-mini-0196:phyxminiproblems-problemphyxmini0196-speedinmeterspersecond, def:physics:phyx-mini-0196:phyxminiproblems-problemphyxmini0196-sirenlistenerdopplersetup}
  The classical one-dimensional Doppler law for a wave travelling from S to L, opposite the positive L-to-S road direction. Both velocities are measured relative to the air, so positive listener velocity raises the encounter rate while positive source velocity lengthens the trailing wavelength. This is a general governing relation and contains no displayed answer value
\end{definition}
\begin{proof}
  This is the declared model definition of Satisfies Leftward Acoustic Doppler Law.
\end{proof}
\begin{definition}[Answer Choice]
  \label{def:physics:phyx-mini-0196:phyxminiproblems-problemphyxmini0196-answerchoice}
  \lean{PhyXMiniProblems.ProblemPhyXMini0196.AnswerChoice}
  Labels of the four answers printed with the problem
\end{definition}
\begin{proof}
  This is the declared model definition of Answer Choice.
\end{proof}
\begin{definition}[hertz]
  \label{def:physics:phyx-mini-0196:phyxminiproblems-problemphyxmini0196-answerchoice-hertz}
  \lean{PhyXMiniProblems.ProblemPhyXMini0196.AnswerChoice.hertz}
  \uses{def:physics:phyx-mini-0196:phyxminiproblems-problemphyxmini0196-answerchoice}
  Whole-hertz frequency printed beside each answer label
\end{definition}
\begin{proof}
  This is the declared model definition of hertz.
\end{proof}
\begin{definition}[recorded Dataset Answer]
  \label{def:physics:phyx-mini-0196:phyxminiproblems-problemphyxmini0196-recordeddatasetanswer}
  \lean{PhyXMiniProblems.ProblemPhyXMini0196.recordedDatasetAnswer}
  \uses{def:physics:phyx-mini-0196:phyxminiproblems-problemphyxmini0196-answerchoice}
  The answer label recorded in the supplied dataset metadata
\end{definition}
\begin{proof}
  This is the declared model definition of recorded Dataset Answer.
\end{proof}
\begin{definition}[Matches Answer Choice]
  \label{def:physics:phyx-mini-0196:phyxminiproblems-problemphyxmini0196-matchesanswerchoice}
  \lean{PhyXMiniProblems.ProblemPhyXMini0196.MatchesAnswerChoice}
  \uses{def:physics:phyx-mini-0196:phyxminiproblems-problemphyxmini0196-acousticfrequency, def:physics:phyx-mini-0196:phyxminiproblems-problemphyxmini0196-frequencyinhertz, def:physics:phyx-mini-0196:phyxminiproblems-problemphyxmini0196-answerchoice}
  Agreement with a displayed whole-hertz answer. A half-hertz tolerance models rounding to the nearest integer without identifying the physical frequency itself with an exact integer
\end{definition}
\begin{proof}
  This is the declared model definition of Matches Answer Choice.
\end{proof}
\begin{lemma}[heard Frequency hertz eq calibrated Ratio]
  \label{lem:physics:phyx-mini-0196:phyxminiproblems-problemphyxmini0196-heardfrequency-hertz-eq-calibratedratio}
  \lean{PhyXMiniProblems.ProblemPhyXMini0196.heardFrequency_hertz_eq_calibratedRatio}
  \uses{def:physics:phyx-mini-0196:phyxminiproblems-problemphyxmini0196-frequencyinhertz, def:physics:phyx-mini-0196:phyxminiproblems-problemphyxmini0196-sirenlistenerdopplersetup, def:physics:phyx-mini-0196:phyxminiproblems-problemphyxmini0196-matchesprimaryfigure, def:physics:phyx-mini-0196:phyxminiproblems-problemphyxmini0196-usespriorsirencalibration, def:physics:phyx-mini-0196:phyxminiproblems-problemphyxmini0196-hassubsonicphysicalparameters, def:physics:phyx-mini-0196:phyxminiproblems-problemphyxmini0196-satisfiesleftwardacousticdopplerlaw}
  Substitution of the figure and calibration readouts into the governing law. The exact ratio is a derived conclusion, not a setup field or premise
\end{lemma}
\begin{proof}
  The conclusion follows from the governing relations and figure readouts named in the statement of heard Frequency hertz eq calibrated Ratio.
\end{proof}
% --- Archon declaration topology end ---
% --- Archon physics formalization source end ---
